\subsection{Direct Preference Optimization}

\textbf{Direct Preference Optimization (DPO)} is a method for aligning language models using preference data, like RLHF, but without training a separate reward model and without running PPO~\cite{rafailov2023direct}.

We start with a normal language model. For a prompt \(x\), the model generates two answers, and humans choose a preferred one. Denote the chosen answer and the rejected answer as \(y^+\) and \(y^-\). Our goal is to fine-tune the model so that
\[
	\pi_\theta (y^+ | x) > \pi_\theta (y^- | x),
\]
where \(\pi_\theta\) is the model being trained.

Then, DPO compares the new model \(\pi_\theta\) to a reference model \(\pi_\text{ref}\), which is usually the fine-tuned model trained from demonstrations in RLHF (\ref{sec:rlhf}).

For one preference pair \((x, y^+, y^-)\), DPO looks at the difference between two log-probability ratios
\[
	\Delta r_\theta =
	\log\frac{\pi_\theta(y^+|x)}{\pi_\text{ref}(y^+|x)} -
	\log\frac{\pi_\theta(y^-|x)}{\pi_\text{ref}(y^-|x)}
\]

This ratio tells whether the model (1) increases the chosen answer and (2) decreases the rejected answer, relative to the reference model. The loss of DPO is
\begin{equation}
	\mathcal{L}_\text{DPO} = - \log \sigma \left( \beta (\Delta r_\theta) \right).
\end{equation}

DPO is simpler and thus replaces PPO in many settings. It collects preference data and fine-tunes directly on chosen-vs-rejected pairs.

The “proximal” in DPO mainly shows up through the objective acting like regularization. Suppose the model tries to make \(y^+\) insanely more likely than the reference model did, then
\[
	\log\frac{\pi_\theta(y^+|x)}{\pi_\text{ref}(y^+|x)}
\]
gets large. As a result, the sigmoid will saturate, and the training gain becomes small. From another perspective, DPO is not simply increasing \(\pi_\theta (y^+|x)\); its objective expects that
\[
	\log\frac{\pi_\theta(y^+|x)}{\pi_\text{ref}(y^+|x)} \quad
	\text{should be larger than} \quad
	\log\frac{\pi_\theta(y^-|x)}{\pi_\text{ref}(y^-|x)}.
\]

So the model is judged by how it changes relative to the reference model, not just by raw probabilities.
